\documentclass[10pt,twocolumn]{article}
\usepackage[a4paper,margin=1.9cm,columnsep=0.7cm]{geometry}
\usepackage{mathptmx}
\usepackage[T1]{fontenc}
\usepackage{microtype}
\usepackage{graphicx}
\usepackage{booktabs}
\usepackage{caption}
\usepackage{authblk}
\usepackage[hidelinks]{hyperref}
\usepackage{xcolor}

\captionsetup{font=small,labelfont=bf,skip=5pt}
\setlength{\parindent}{1.1em}
\definecolor{mono}{HTML}{2F6F9F}
\definecolor{dia}{HTML}{C25A3C}

\title{\vspace{-1.2em}\bfseries Monologue or Dialogue?\\[2pt]
\large Measuring Conversational Structure in the Reasoning Traces of Large Language Models}
\author[ ]{Oskar Engl}
\affil[ ]{\small Johannes Kepler University Linz, Austria \quad\textit{oskar.engl@ik.me}}
\date{\small Condensed version --- July 2026}

\begin{document}
\twocolumn[
\begin{@twocolumnfalse}
\maketitle
\begin{abstract}
\noindent
A reasoning model such as DeepSeek-R1 emits a long ``thinking'' trace before it answers. That
trace can read as one mind reasoning step by step, or as several voices trading proposals and
objections. We ask which description the \emph{text} fits when the question is put to measurement.
Because no reasoning trace comes labelled with its true voice structure, we build a $2\times2$
grid of corpora whose structure \emph{is} known --- human and AI, monologue and dialogue --- and
admit a detection method to the traces only once it can tell monologue from dialogue on the four
known corpora. Three methods qualify, each reading a different signal: a document-level
author-count detector, a sentence-level style-switch detector, and a dialogue-act detector that
measures adjacency-pair completion. Applied to the traces, \emph{two of the three} place them on
the dialogue side, past the dialogue pole --- the sentence-level style switching and the
adjacency-pair completion --- while the coarse document-level author-count detector leans
monologue, and only marginally. Read together, R1 traces carry dialogue-like structure both in
their sentence-to-sentence style and in their question-and-answer organisation; only the bluntest
measure, whole-document authorship, fails to resolve them into distinct voices.
\medskip
\end{abstract}
\end{@twocolumnfalse}
]

\begin{figure*}[t]
\centering
\includegraphics[width=0.96\textwidth]{cs_fig_grid.pdf}
\caption{\textbf{The design.} Left to right: four corpora of known voice structure feed a
\emph{validity gate} that each detection method must pass --- separating monologue from dialogue
on all four known cells. All three methods qualify and are then applied to the DeepSeek-R1 traces
(\texttt{ai-reason}), which they place on the monologue--dialogue axis anchored by the two
same-task AI cells; the human cells validate only.}
\label{fig:grid}
\end{figure*}

\section{The question}
Reasoning-oriented training is now standard in frontier models: given a problem, a model like
DeepSeek-R1 \cite{guo2025} first generates an extended stretch of intermediate tokens
\cite{kambhampati2025} --- a reasoning trace --- and only then commits to an answer. These traces
have a curious property. Each is produced by a single model in one pass, yet a trace often contains
questions and their answers, objections and concessions, and abrupt reversals
(``Wait, that's not right''). In human writing such features can mark a genuine exchange, but they
can equally mark one vivid inner voice. Hence a question easy to state and less easy to settle:
\emph{does a reasoning trace have the textual structure of a monologue or of a dialogue?}

The question has a live theoretical stake. Kim et al.~\cite{kim2026} argue that reasoning-trained
models improve partly because they simulate a ``society of thought'' --- distinct internal
perspectives that propose, criticise, and reconcile --- and older ideas echo it: reasoning as at
root argumentative and social \cite{mercier2017}, and the mind as many interacting agents
\cite{minsky1986}. Our premise is narrower and testable: voice structure leaves marks in the
writing itself. If distinct perspectives take turns producing the text, they should also differ in
the small, largely involuntary habits --- function-word use, sentence rhythm, dialogue-act moves
--- that authorship and conversation analysis use to tell writers and turns apart. We therefore
treat the question operationally: a monologue should read as one author, a dialogue as several
voices exchanging turns, and we ask which the traces resemble.

\section{The design: a calibration grid}
No corpus of reasoning traces is labelled with its true voice structure, which is what makes the
question hard to settle directly. We address this with the grid in Figure~\ref{fig:grid}. Four
corpora of \emph{known} structure vary producer (human vs.\ AI) and structure (monologue vs.\
dialogue): human monologue (think-aloud protocols \cite{wurgaft2025}), human dialogue (group
deliberation, DeliData \cite{karadzhov2023}), AI monologue (one Qwen-2.5-7B instance solving a
puzzle alone), and AI dialogue (four stateless Qwen-2.5-7B instances, each given a fixed role,
debating). The fifth cell --- native DeepSeek-R1 traces --- is the target.

The grid works as a \emph{gate}: a method may report on the target only after it correctly
separates monologue from dialogue on all four known cells. A detector that cannot tell voices apart
in a genuine four-way debate is in no position to rule them out in a trace; one that reports several
voices in a lone think-aloud protocol is in no position to affirm them there. This is a validity
check the method must earn, not an assumption we grant it.

One design choice keeps the contrast clean: we hold the task fixed wherever possible.
\texttt{ai-mono}, \texttt{ai-dia}, \texttt{ai-reason} and \texttt{hum-mono} all use the same puzzle,
Game of 24 \cite{yao2023}, so a method reading a monologue as one voice cannot merely be reacting
to subject matter. (\texttt{hum-dia} uses the Wason task, the closest human group-deliberation data
available; its role is only to confirm that a method detects human dialogue at all, so it never
anchors the placement.) Speaker-identity tokens --- pseudonyms, role titles, letter labels --- were
stripped from both dialogue cells before analysis, so dialogue detection cannot be an artefact of
speaker cues rather than interactional structure.

\section{Three methods, three signals}
The methods read deliberately different signals, so that agreement among them carries more weight
than any one alone.

\textbf{Method 1 --- document-level author count} \cite{nath2019}. Clusters a document's
paragraph-sized windows on their 50-most-frequent-token profiles and reports the number of
clusters as an author count, reduced to a binary single/multi call. \emph{Gate:} balanced accuracy
across the four cells must beat the $0.50$ chance baseline.

\textbf{Method 2 --- sentence-level style switch} \cite{schmidt2025}. A pretrained sequential
sentence-pair classifier emits, for each adjacent sentence pair, the probability that the two were
written by different authors, using content-independent style embeddings \cite{patel2025}. Because
raw scores depend on length, each document's score is turned into a length-adjusted \emph{lean}
--- its distance to the dialogue anchor's fitted length curve minus its distance to the monologue
anchor's. \emph{Gate:} mean lean must be higher on the dialogue cells than the monologue cells,
with Cliff's $\delta\ge0.33$.

\textbf{Method 3 --- dialogue-act structure} \cite{duran2022,schegloff1973}. An LLM tags each
utterance with a CAMS dialogue act and measures adjacency-pair completion: when a first pair-part
(a question, request or offer) occurs, is it met by a valid second pair-part (an answer, a
confirmation, an acceptance) in the very next utterance --- the mark of conversation a lone writer
has little reason to produce. \emph{Gate:} the completion rate must separate the dialogue cells from
the monologue cells, with Cliff's $\delta\ge0.33$.

\section{Results}
\textbf{Method 1 passes, but softly.} Balanced accuracy is $0.6088$ (bootstrap 95\% CI
$[0.596,0.621]$), clearing the chance baseline. But the pass leans on near-perfect single-author
recall ($0.975$ on \texttt{hum-mono}, $0.998$ on \texttt{ai-mono}) averaged against weak
multi-author recall: $0.423$ on \texttt{ai-dia} and a complete miss on \texttt{hum-dia} ($0.000$).
The detector flags dialogue where it finds it but essentially never rules single authorship
\emph{out}. A monologue-leaning placement under this method should thus be read as the
\emph{absence} of a detected authorship signal, not positive evidence of one voice.

\textbf{Method 2 passes cleanly.} Mean lean is higher on the dialogue cells than the monologue cells
($0.0125$ vs.\ $-0.0105$) with Cliff's $\delta=0.6211$, well above the $0.33$ floor. The traces
themselves ($0.0433$) sit above \emph{both} AI anchors.

\textbf{Method 3 passes.} Its adjacency-pair completion rate separates dialogue from monologue in
the expected direction, with Cliff's $\delta=0.4717$, above the $0.33$ floor: within each producer
the dialogue cell completes pairs more often than the monologue cell (human $0.235$ vs.\ $0.096$;
AI $0.025$ vs.\ $0.003$). Its two descriptive signals --- the distribution of dialogue-act tags and
their transitions --- do not separate the AI poles, because every cell is dominated by the
\texttt{inform} act; the informative signal is completion, on which the method rests.

\begin{figure*}[t]
\centering
\includegraphics[width=0.83\textwidth]{cs_fig_split.pdf}
\caption{\textbf{Where the three methods place the traces.} Each method's placement of the target on
the monologue--dialogue axis (poles at $0$ and $1$). The shaded band is that method's margin, within
which no lean is called. Method~1's marker sits at the edge of a \emph{wide} margin, just onto the
monologue side --- a marginal call. Methods~2 and~3 both land well past the dialogue pole. Two of
the three point to dialogue.}
\label{fig:split}
\end{figure*}

\textbf{Two of the three point to dialogue} (Table~\ref{tab:verdict}, Figure~\ref{fig:split}).
Method~2, reading sentence-to-sentence style, places the traces firmly on the \emph{dialogue} side,
past the dialogue anchor itself. Method~3, reading the completion of question-and-answer pairs,
places them past the dialogue anchor as well: the trace completion rate ($0.035$) sits closer to
the dialogue anchor \texttt{ai-dia} ($0.025$) than to the monologue anchor \texttt{ai-mono}
($0.003$). Only Method~1, the coarsest measure, leans \emph{monologue}, and it does so marginally
--- its marker sits at the very edge of a wide no-decision band. The two methods that read
fine-grained interactional structure agree; the disagreement is confined to the one method that
reads only coarse, document-level authorship.

\begin{table}[t]
\centering
\small
\caption{Cross-method verdicts on the reasoning traces. This is the answer to the research
question.}
\label{tab:verdict}
\begin{tabular}{@{}lcc@{}}
\toprule
Method & Passed gate? & Verdict on \texttt{ai-reason} \\
\midrule
1: author count (doc.) & pass & \textcolor{mono}{monologue (marginal)} \\
2: style switch (sent.) & pass & \textcolor{dia}{dialogue} \\
3: adjacency pairs & pass & \textcolor{dia}{dialogue} \\
\bottomrule
\end{tabular}
\end{table}

\section{What it means}
The claim this study licenses is about \emph{textual surface structure}, not hidden minds: whether
the traces exhibit the style-switching and adjacency-pair regularities of multi-party dialogue. On
that question the two methods that measure fine-grained interactional structure agree: the traces
switch style from sentence to sentence like a dialogue, and they complete question-and-answer pairs
at a rate above the dialogue anchor. The lone monologue call comes from the coarsest measure,
whole-document author clustering --- and that method is a conservative detector whose monologue
verdict means ``no authorship signal found,'' not ``positively one voice.'' The weight of the
evidence therefore tilts toward dialogue.

This aligns the outside view with the inside one. Kim et al.~\cite{kim2026} infer distinct internal
perspectives from interpretability and trace labelling; the surface text measured here carries the
matching marks --- style that shifts across the trace and first pair-parts met by their second
parts --- which is the measurable surface counterpart of the internal multi-agent dynamics they
report. A dialogue-leaning reading does not show that a model contains multiple minds, any more than
a novelist's dialogue does; it shows that the trace is organised as conversation, itself a
substantive, measurable fact.

Two qualifications temper the reading, and both concern how a distance to a pole is obtained rather
than its direction. Method~2's monologue distance is \emph{extrapolated}: the traces run far longer
than the monologue anchor, so what the method measures directly is that they track the
\emph{dialogue} anchor closely across the shared length range, while the monologue comparison
extends a length curve past the data that fit it. And Method~3's verdict rests specifically on the
adjacency-pair completion signal; its coarser tag profiles cannot separate the AI poles at all.
Neither qualification reverses a direction --- both dialogue-ward placements survive them --- but
each bounds how far past the pole the placement should be read.

\section{Limitations}
The target is one model (DeepSeek-R1) on one puzzle family, so the verdict is about this cell, not
reasoning traces at large. The target is also a \emph{different} base model than the anchors
(Qwen-2.5-7B), so any measured distance carries a shared model offset; the placement is best read
as ``closer to which pole,'' not as an absolute distance. Method~3's dialogue-act labels are
produced by an automatic tagger, and Method~2's monologue comparison rests on extrapolation beyond
the fitted length range. These are the reasons the validity gate matters: a method that survives it
has at least shown it can tell the two structures apart on data where the answer is known.

\begin{thebibliography}{99}
\small
\setlength{\itemsep}{0pt}\setlength{\parskip}{0pt}
\bibitem{guo2025} D.\ Guo et al. DeepSeek-R1 incentivizes reasoning in LLMs through reinforcement
learning. \textit{Nature} 645:633--638, 2025.
\bibitem{kambhampati2025} S.\ Kambhampati et al. Stop anthropomorphizing intermediate tokens as
reasoning/thinking traces! \textit{arXiv:2504.09762}, 2025.
\bibitem{kim2026} J.\ Kim, S.\ Lai, N.\ Scherrer, B.\ Ag\"uera y Arcas, J.\ Evans. Reasoning models
generate societies of thought. \textit{arXiv:2601.10825}, 2026.
\bibitem{mercier2017} H.\ Mercier, D.\ Sperber. \textit{The Enigma of Reason.} Harvard Univ.\
Press, 2017.
\bibitem{minsky1986} M.\ Minsky. \textit{The Society of Mind.} Simon \& Schuster, 1986.
\bibitem{wurgaft2025} D.\ Wurgaft et al. Scaling up the think-aloud method. \textit{Proc.\ CogSci},
2025.
\bibitem{karadzhov2023} G.\ Karadzhov, T.\ Stafford, A.\ Vlachos. DeliData: a dataset for
deliberation in multi-party problem solving. \textit{Proc.\ ACM HCI} 7(CSCW2), 2023.
\bibitem{yao2023} S.\ Yao et al. Tree of thoughts: deliberate problem solving with LLMs.
\textit{NeurIPS} 36, 2023.
\bibitem{nath2019} S.\ Nath. Style change detection by threshold-based and window-merge clustering.
\textit{CLEF Working Notes}, 2019.
\bibitem{schmidt2025} G.\ Schmidt et al. Style change detection using a sequential sentence-pair
classifier. \textit{CLEF Working Notes}, 2025. \textit{arXiv:2508.00675}.
\bibitem{patel2025} A.\ Patel et al. StyleDistance: content-independent style embeddings with
synthetic parallel examples. \textit{NAACL}, 2025.
\bibitem{duran2022} N.\ Duran, S.\ Battle, J.\ Smith. Inter-annotator agreement using the
Conversation Analysis Modelling Schema. \textit{Comm.\ Methods \& Measures} 16(3):182--214, 2022.
\bibitem{schegloff1973} E.\ Schegloff, H.\ Sacks. Opening up closings. \textit{Semiotica}
8(4):289--327, 1973.
\end{thebibliography}

\end{document}
